\icode{E-MD-CA-L3-261113-001}
如图所示，传送带与水平面夹角 \(\theta =\) \,\qty{37}{\degree}，以恒定速率 \(v =\) \,\qty{10}{m/s}沿顺时针方向转动。	现在传送带上端 \(A\) 处无初速度地放一质量 \(m =\) \,\qty{1}{kg}的小煤块（可视为质点，忽略滑动过程中的质量损失），小煤块与传送带间的动摩擦因数 \(\mu = 0.5\)，已知传送带上 \(A\) 到 \(B\) 的长度 \(L =\) \,\qty{16}{m}。取\(\sin\) \,\qty{37}{\degree}\(= 0.6\)，\(\cos\) \,\qty{37}{\degree}\(= 0.8\)，重力加速度 \(g =\) \,\qty{10}{m/s^2}。求：
\begin{enumerate}
	\item 小煤块受到的摩擦力大小；
	\item 小煤块刚开始运动时的加速度大小；
	\item 小煤块从 \(A\) 运动到 \(B\) 的时间。
\end{enumerate}

\begin{tikzpicture}[>=stealth,font=\scriptsize]
	\begin{scope}[rotate=37]
		\coordinate (o) at(0,0);
		\coordinate (length) at(4,0);
		\coordinate (square) at(3.7,.3);
		\draw[gray] (0,0) circle(.3) (length) circle(.3);
		\draw[fill=orange!50] (square) rectangle($(square)+(.3,.3)$);
		\draw ($(o)+(0,.3)$)--($(length)+(0,.3)$) arc(90:-90:.3)--($(length)+(0,-.3)$)--($(o)+(0,-.3)$) arc(270:90:.3);
		\fill (o) circle(.07) (length) circle(.07);
		\draw[->] (1.5,.15)--node[below,rotate=37]{\(v\)}(2.5,.15);
		\node at(4,.3)[above right]{\(A\)};
		\node at(0,.3)[above left]{\(B\)};
	\end{scope}
	\draw[dashed,cyan] (0,-.3)--(2,-.3);
	\draw[gray] ({.3*tan(18.5)+.4},-.3) arc(0:37:.4);
	\node at(.6,-.12){\(\theta\)};
\end{tikzpicture}
\begin{answer}
    
\end{answer}